\chapter{Menemukan dan Memampatkan Dinamika}
\label{ch:penemuan}

Makalah SINDy \citep{brunton2016} sudah saya simpan sejak musim panas 2024, sebelum kuliah pertama
dimulai, dan buku \citet{brunton2022} menyusul ke folder bacaan tidak lama kemudian. Bab ini membahas
gagasan-gagasan dari buku itu yang menyambung dengan banyak kuliah: PCA dari kuliah unsupervised dan
computer vision, nilai eigen dari fisika komputasi, lasso dari dasar-dasar machine learning, dan sistem
dinamis dari kuliah kendali.

Model pengganti di \cref{ch:operator} umumnya kotak hitam yang besar sekali. Bab ini menempuh arah
sebaliknya. Pertanyaannya bukan ``bisakah jaringan meniru simulasi?'', melainkan ``berapa sedikit yang
sebenarnya kita butuhkan untuk menjelaskan dinamika ini?''. Jawabannya sering mengejutkan: medan aliran
dengan jutaan derajat kebebasan kadang bisa dijelaskan oleh beberapa mode saja, dan dinamikanya oleh
persamaan dengan beberapa suku.

\section{Proper orthogonal decomposition}

Simpan $m$ potret medan aliran, masing-masing dengan $n$ nilai, sebagai kolom-kolom matriks
$\mX\in\R^{n\times m}$ setelah rata-ratanya dikurangkan. Dekomposisi nilai singular memberi
\begin{equation}
\mX=\bm U\bm\Sigma\bm V^\top,
\end{equation}
dengan kolom-kolom $\bm U$ sebagai mode ruang, nilai singular $\sigma_k$ di diagonal $\bm\Sigma$, dan
baris-baris $\bm V^\top$ sebagai evolusi waktu setiap mode. Dalam mekanika fluida, prosedur ini disebut
\istilah{proper orthogonal decomposition} \citep{berkooz1993,taira2017}. Ia tidak lain adalah PCA dari
\cref{ch:unsup}, diterapkan pada potret-potret medan.

Mengapa mode-mode ini istimewa? Teorema Eckart--Young menyatakan bahwa di antara semua matriks berpangkat
$r$, potongan SVD yang hanya menyimpan $r$ mode pertama adalah yang paling dekat ke $\mX$, dengan galat
kuadrat sama dengan jumlah kuadrat nilai singular yang dibuang, $\sum_{k>r}\sigma_k^2$. Kalau medan berupa
kecepatan, $\sigma_k^2$ sebanding dengan energi kinetik yang ditangkap mode ke-$k$, sehingga fraksi
$\sum_{k\le r}\sigma_k^2/\sum_k\sigma_k^2$ dibaca sebagai fraksi energi yang ditangkap. Untuk aliran
melewati silinder pada bilangan Reynolds rendah, beberapa mode pertama datang berpasangan dan bersama-sama
menangkap hampir seluruh energi fluktuasi, karena pelepasan pusaran adalah gelombang yang berjalan,
dan gelombang berjalan butuh dua mode yang bergeser seperempat periode.

Mode POD bisa dipakai untuk membangun \istilah{reduced-order model}. Tuliskan
$\vu(\vx,t)\approx\bar\vu(\vx)+\sum_{k=1}^ra_k(t)\bm\varphi_k(\vx)$, masukkan ke persamaan Navier--Stokes,
lalu proyeksikan ke setiap mode. Proyeksi Galerkin ini menghasilkan sistem ODE kecil untuk koefisien
$a_k(t)$, dengan suku linear dari viskositas dan suku kuadratik dari konveksi. Model seperti ini bisa
ribuan kali lebih cepat, tetapi sering tidak stabil untuk waktu panjang, karena mode-mode kecil yang
dibuang adalah tempat energi seharusnya dibuang, sama seperti pusaran kecil dalam kaskade turbulensi
(\cref{ch:numerik}). Model penutup untuk mode yang hilang, atau model deret seperti LSTM pada koefisien
$a_k$ (\cref{ch:lstm}), dipakai untuk memperbaikinya.

\section{Dynamic mode decomposition}

POD memilih mode menurut energi, tanpa peduli dinamikanya. \istilah{Dynamic mode decomposition}
\citep{schmid2010} memilih mode menurut dinamika. Misalkan potret-potret berurutan dengan jarak waktu
$\Delta t$, dan susun dua matriks yang bergeser satu langkah, $\mX=[\vx_1,\dots,\vx_{m-1}]$ dan
$\mX'=[\vx_2,\dots,\vx_m]$. Kita mencari operator linear $\mA$ sehingga $\mX'\approx\mA\mX$, yaitu
$\mA=\mX'\mX^+$ dengan pseudo-invers $\mX^+$. Untuk medan besar, $\mA$ berukuran $n\times n$ dan
terlalu besar untuk dibentuk, sehingga kita memproyeksikannya ke $r$ mode POD pertama. Dengan
$\mX\approx\bm U_r\bm\Sigma_r\bm V_r^\top$, diperoleh
\begin{equation}
\tilde{\mA}=\bm U_r^\top\mX'\bm V_r\bm\Sigma_r^{-1},
\end{equation}
matriks kecil $r\times r$ yang nilai eigennya sama dengan nilai eigen utama $\mA$. Setiap nilai eigen
$\lambda_k$ membawa satu mode yang berkembang seperti $\lambda_k^t$ per langkah. Dalam waktu kontinu,
$\omega_k=\ln\lambda_k/\Delta t$, dengan bagian real sebagai laju tumbuh atau luruh dan bagian imajiner
sebagai frekuensi. Sebagai contoh, $\lambda=0{,}99\,e^{0{,}1i}$ dengan $\Delta t=0{,}01$ detik memberi
laju luruh $\ln0{,}99/0{,}01\approx-1{,}0$ per detik dan frekuensi sudut 10 radian per detik.

DMD seperti mendengarkan sebuah orkestra lalu menuliskan partitur: nada mana yang dimainkan, seberapa
keras, dan apakah ia menguat atau memudar. POD, sebaliknya, hanya mencatat alat musik mana yang paling
banyak bersuara.

\section{Koopman: linear di ruang yang lebih besar}

DMD mendekati dinamika dengan operator linear, padahal kebanyakan sistem nonlinear. Teori operator
Koopman \citep{koopman1931} memberi pembenarannya. Alih-alih melacak keadaan $\vx$, kita melacak fungsi-fungsi
pengamatan $g(\vx)$. Operator Koopman memajukan pengamatan dalam waktu, dan operator itu selalu linear,
walau dinamikanya nonlinear, dengan harga bahwa ruang fungsi pengamatan umumnya berdimensi tak hingga.

Contoh dari \citet{brunton2022} memperlihatkan kapan ruang itu ternyata berhingga. Ambil sistem
$\dot x=\mu x$ dan $\dot y=\lambda(y-x^2)$, yang nonlinear karena suku $x^2$. Tambahkan pengamatan baru
$z=x^2$. Turunannya $\dot z=2x\dot x=2\mu x^2=2\mu z$. Dalam koordinat $(x,y,z)$, sistemnya menjadi
\begin{equation}
\dot x=\mu x,\qquad \dot y=\lambda y-\lambda z,\qquad \dot z=2\mu z,
\end{equation}
sepenuhnya linear. Dalam praktik, pengamatan yang membuat dinamika linear jarang bisa ditebak, sehingga
\citet{lusch2018} melatih autoencoder untuk menemukannya: encoder memetakan keadaan ke koordinat di mana
dinamikanya adalah perkalian matriks, dan decoder mengembalikannya. Kalau berhasil, semua alat kendali
linear dari \cref{ch:kendali} bisa dipakai pada sistem nonlinear.

\section{Menemukan persamaan: SINDy}

Kadang kita ingin lebih dari model yang meniru: kita ingin persamaannya. \istilah{Sparse identification
of nonlinear dynamics} \citep{brunton2016} berangkat dari pengamatan bahwa kebanyakan persamaan fisika
hanya punya sedikit suku. Susun pustaka fungsi kandidat dari data, misalnya
$\bm\Theta(\mX)=[\bm1,\ \vx,\ \vx^2,\ \vx^3,\ \sin\vx,\dots]$, dan cari koefisien jarang $\bm\Xi$ sehingga
\begin{equation}
\dot{\mX}\approx\bm\Theta(\mX)\,\bm\Xi.
\end{equation}
Ini masalah regresi linear dengan prior jarang, persis lasso dari \cref{ch:data}. Algoritme yang dipakai
dalam makalah aslinya lebih sederhana: selesaikan kuadrat terkecil, nolkan koefisien yang kecil di bawah
ambang, lalu ulangi kuadrat terkecil hanya dengan suku yang tersisa.

Contoh kecil memperlihatkan logikanya. Misalkan data berasal dari $\dot x=-2x$, dan pustakanya
$[1,x,x^2]$. Pada data tanpa noise, $\dot x$ di setiap titik tepat $-2$ kali $x$, sehingga kuadrat
terkecil memberi koefisien $(0,-2,0)$ dan pengambangan tidak mengubah apa pun. Dengan noise, kuadrat
terkecil biasa memberi koefisien kecil pada $1$ dan $x^2$ yang mencoba menjelaskan noise. Pengambangan
membuangnya, dan regresi ulang hanya dengan $x$ memulihkan $-2$ dengan baik. Hasilnya bukan jaringan
dengan ribuan bobot, melainkan persamaan yang bisa dibaca, diperiksa, dan diekstrapolasi. Pada data
sistem Lorenz yang chaos, SINDy memulihkan ketiga persamaannya hanya dari lintasan yang diukur.

SINDy punya tiga syarat yang sering menjadi titik lemah. Turunan waktu harus diperkirakan dari data,
dan diferensiasi numerik memperbesar noise. Pustakanya harus memuat suku yang benar. Dan koordinatnya
harus tepat, karena dalam koordinat yang salah, dinamika sederhana bisa tampak rumit. Untuk masalah
terakhir, \citet{champion2019} melatih autoencoder bersama SINDy, sehingga koordinat dan persamaan
ditemukan bersamaan. PDE-FIND \citep{rudy2017} memperluas gagasan ini ke persamaan diferensial parsial
dengan memasukkan turunan ruang ke dalam pustaka.

\section{Memilih alat}

Ketiga keluarga ini mengisi tempat yang berbeda. POD dan model tereduksi cocok ketika kita punya
persamaan dan ingin mempercepatnya. DMD dan Koopman cocok untuk menganalisis dan meramalkan dinamika
yang hampir periodik atau berosilasi, dan untuk kendali. SINDy cocok ketika persamaannya belum diketahui
dan datanya cukup bersih. Semuanya menghasilkan model yang jauh lebih kecil dan lebih mudah ditafsirkan
daripada jaringan dalam, dan karena itulah mereka masih dipakai berdampingan dengan neural operator di
\cref{ch:operator}, sering sebagai lapisan di dalamnya.

\sumberbab{Bab ini dirangkai dari kuliah Machine Learning: Unsupervised Techniques (PCA), Control Systems
(sistem linear dan nilai eigen), Computational Physics I (SVD dan masalah nilai eigen), dan bacaan selama
tugas akhir. Rujukan utama: \citet{brunton2022}, \citet{taira2017}, \citet{schmid2010}, dan
\citet{brunton2016}.}
